% ============================================================
% Capitolo 7 — Conclusioni e lavoro futuro (traduzione italiana di
% chapters/ch07_conclusions.tex; label invariati)
% ============================================================
\chapter{Conclusioni e lavoro futuro}
\label{ch:conclusions}

\section{Conclusioni}
\label{sec:conclusions}

Questa tesi ha chiesto se ordinare i task di addestramento dal facile
al difficile produca un comportamento di guida multi-agente
misurabilmente diverso rispetto al presentare tutte le difficoltà
insieme, a parità di budget, in un ambiente CARLA a mondo condiviso
addestrato con MAPPO. La risposta è sì --- ma non lungo l'asse su cui i
curricula vengono di solito venduti.

\begin{itemize}
  \item \textbf{L'ordinamento non ha accelerato l'apprendimento.} Tutte
        e quattro le celle della campagna hanno raggiunto le stesse
        soglie iniziali di competenza agli stessi conteggi di passi,
        plausibilmente perché il reward shaping denso fornisce già
        segnale di apprendimento sui task difficili
        (\cref{sec:sample-efficiency}, \cref{sec:ordering}).
  \item \textbf{L'ordinamento ha selezionato ciò che la policy è
        diventata.} Sotto il critic MLP, il curriculum ha prodotto un
        guidatore cauto e incline allo stallo e il batch uno veloce e
        incline alle collisioni; sul criterio di successo rigoroso lo
        scambio si è risolto a favore del curriculum di $+18.5$\,pp in
        valutazione, coerentemente in tutti gli scenari
        (\cref{sec:final-performance}).
  \item \textbf{L'effetto è condizionato dalla stima del valore.} Sotto
        un critic GNN l'effetto di regime in-domain è svanito
        ($-0.1$\,pp) mentre l'architettura ha recuperato la cella batch
        e rimodellato la cella curriculum --- una forte interazione
        architettura$\times$regime (\cref{sec:arch-axis}).
  \item \textbf{Il beneficio portabile è la generalizzazione.} L'unica
        componente del vantaggio del curriculum replicata in entrambe
        le architetture è il margine su una città mai vista in
        addestramento ($+19.0$ e $+10.0$\,pp su Town05), contro
        l'intuizione copertura-implica-robustezza
        (\cref{sec:generalization}).
  \item \textbf{Specificità di classe e limiti della piattaforma.} I
        pedoni sono risultati in larga parte insensibili al regime,
        vincolati dalla fattibilità dei percorsi più che dalla qualità
        della policy (\cref{sec:ped-limits}); e ogni cella ha mostrato
        erosione tardiva entro la run a livello di task fissato,
        indipendente dal regime (\cref{sec:behavior}).
\end{itemize}

Oltre ai risultati, la tesi contribuisce l'apparato che li ha resi
misurabili: un framework di confronto a parità di budget con
configurazioni appaiate per seed e byte-identiche; un protocollo di
sviluppo guidato da gate con un registro esplicito di modifiche
promosse, rigettate e consapevolmente mantenute; un resoconto
documentato di come i difetti di validità dell'ambiente possano
mascherarsi da effetti di policy, con le loro firme di rilevamento; e
un metro di varianza da repliche a seed identico contro cui ogni
affermazione viene letta. In un contesto in cui un singolo bug del
simulatore ha spostato la metrica principale di cinquanta punti
percentuali, questa disciplina di misura --- non un singolo numero ---
è ciò che rende il confronto significativo.

\section{Lavoro futuro}
\label{sec:future}

In ordine di quanto direttamente ogni voce estende questa tesi:

\begin{enumerate}
  \item \textbf{Replicazione dei seed.} Il design è stato tagliato da
        tre seed appaiati per cella a uno per ragioni di calcolo
        (\cref{sec:campaign}). Rieseguire la matrice $2\times2$ a
        $n{=}3$ trasformerebbe il caso di studio appaiato in stime di
        distribuzione e metterebbe alla prova direttamente l'entità
        dell'interazione.
  \item \textbf{Spiegare l'erosione.} La degradazione entro la run
        controllata per livello (\cref{tab:quartiles}) invita ad
        ablazioni mirate: congelare una classe a metà addestramento per
        isolare la co-adattazione, e variare l'ancora dell'annealing di
        entropia per isolare lo schedule di esplorazione.
  \item \textbf{Generazione dei percorsi pedonali.} Con più della metà
        dei percorsi pedonali sotto target su Town03
        (\cref{sec:ped-limits}), un planner consapevole della
        fattibilità (esplorazione più profonda dei rami, target
        adattivi per regione di spawn) è la correzione d'ambiente a
        maggior leva, e affinerebbe ogni confronto lato pedoni.
  \item \textbf{Estensioni architetturali sul critic.}
        L'infrastruttura espone già encoder basati su attenzione (GAT)
        e la normalizzazione del valore PopArt dietro flag
        (\cref{sec:critic-encoders}); vista la forza dell'interazione
        con l'encoder del critic, queste sono celle successive naturali
        più che aggiunte esotiche.
  \item \textbf{Curricula adattivi.} Il gate di competenza è un
        sequenziatore fisso a due soglie; curricula basati sulla
        prestazione che adattano le quote di livello in continuo (per
        es.\ campionamento proporzionale al progresso di apprendimento)
        metterebbero alla prova se l'effetto di selezione
        dell'equilibrio sopravviva a un ordinamento a grana più fine.
  \item \textbf{Co-design reward--curriculum.} Il tuning della reward
        guidato da gate è girato per lo più su configurazioni bloccate
        su easy (\cref{sec:threats}); ricalibrare i termini chiave di
        shaping sul mix completo quantificherebbe quanto dell'effetto
        di regime è mediato dalla calibrazione della reward.
  \item \textbf{Scalare la cornice.} La copertura parallela dei task è
        stata fissata a $3{+}3$ agenti; scalare il mondo condiviso
        verso decine di agenti --- dove la struttura relazionale del
        critic GNN dovrebbe contare di più --- trasforma la cornice di
        coesistenza stessa in una variabile sperimentale.
\end{enumerate}
